% GENERATED FILE. Edit canonical lessons, then run npm run generate:books.
% canonical-source-hash: fnv1a64:a9baeacec941aadc
% canonical-lessons: HI-C125-shant, HI-C125-ghabrahat, HI-C125-bechain, HI-C125-mannahin, HI-C125-maza

\chapter{Calm and Restless}
\label{ch:hi-calm-and-restless}
\hlchaptermodality{125}

\begin{chapteropening}
\textbf{By the end of this chapter:} \emph{I can say I am calm, nervous or restless, that I don't feel like it, and that I enjoyed something.}

Complete HI-C125-maza: moods, good and uneasy.
\end{chapteropening}

\section[\texorpdfstring{\hi{शांत} (śānt)}{शांत (śānt)}]{\hi{शांत} (śānt) — calm, peaceful}

\label{lesson:HI-C125-shant}



\begin{quote}
Before the new one: say the Hindi for sweat, then the Hindi for a headache.
\end{quote}

\subsection*{You'll want to know: \hi{शांत}}
\textbf{\hi{शांत}} (\emph{śānt}) — \textquotedblleft{}calm, peaceful\textquotedblright{}, of a person or a place.

\textbf{\hi{शांत} \hi{रहिए}} — \textquotedblleft{}please stay calm\textquotedblright{}.

\begin{cousinweb}
\textbf{\hi{शांत}} (\emph{śānt}) is Sanskrit \emph{śānta}, \textquotedblleft{}calm\textquotedblright{}, taken over whole as a learned loan, a tatsama (CDIAL 12390).
\end{cousinweb}

\begin{grammarlens}[title={What you've built}]
Calm.
\end{grammarlens}

\subsection*{Your turn}
\begin{itemize}
\raggedright
  \item \emph{Say it:} \emph{śānt}
  \item \emph{Say it:} \emph{śānt}, once more
  \item \emph{Say it:} say \emph{śānt rahie}
  \item \emph{Recall:} say the Hindi for sweat, then the Hindi for a headache, then say \emph{śānt} again
\end{itemize}

\begin{tcolorbox}[breakable,colback=teal!4,colframe=teal!35!black,title={Before you move on}]
What does \hi{शांत} mean? (\textquotedblleft{}Calm, peaceful\textquotedblright{}.) Say it once more.
\end{tcolorbox}
\section[\texorpdfstring{\hi{घबराहट} (ghabrāhaṭ)}{घबराहट (ghabrāhaṭ)}]{\hi{घबराहट} (ghabrāhaṭ) — nervousness, panic}

\label{lesson:HI-C125-ghabrahat}



\begin{quote}
Before the new one: say the Hindi for a headache, then the Hindi for calm.
\end{quote}

\subsection*{You'll want to know: \hi{घबराहट}}
\textbf{\hi{घबराहट}} (\emph{ghabrāhaṭ}) — \textquotedblleft{}nervousness\textquotedblright{}, a flutter of panic.

\textbf{\hi{मुझे} \hi{घबराहट} \hi{हो} \hi{रही} \hi{है}} — \textquotedblleft{}I'm feeling nervous\textquotedblright{}.

\begin{cousinweb}
\textbf{\hi{घबराहट}} (\emph{ghabrāhaṭ}) is the noun of \emph{ghabrānā}, \textquotedblleft{}to panic\textquotedblright{}. CDIAL lists Hindi \emph{ghābrā}, \textquotedblleft{}startled\textquotedblright{}, under a reconstructed \emph{*ghabbhara}, \textquotedblleft{}confused\textquotedblright{} (CDIAL 4427).
\end{cousinweb}

\begin{grammarlens}[title={What you've built}]
The opposite of calm.
\end{grammarlens}

\subsection*{Your turn}
\begin{itemize}
\raggedright
  \item \emph{Say it:} \emph{ghabrāhaṭ}
  \item \emph{Say it:} \emph{ghabrāhaṭ}, once more
  \item \emph{Say it:} say \emph{mujhe ghabrāhaṭ ho rahī hai}
  \item \emph{Recall:} say the Hindi for a headache, then the Hindi for calm, then say \emph{ghabrāhaṭ} again
\end{itemize}

\begin{tcolorbox}[breakable,colback=teal!4,colframe=teal!35!black,title={Before you move on}]
What does \hi{घबराहट} mean? (\textquotedblleft{}Nervousness, panic\textquotedblright{}.) Say it once more.
\end{tcolorbox}
\section[\texorpdfstring{\hi{बेचैन} (becain)}{बेचैन (becain)}]{\hi{बेचैन} (becain) — restless, uneasy}

\label{lesson:HI-C125-bechain}



\begin{quote}
Before the new one: say the Hindi for calm, then the Hindi for nervousness.
\end{quote}

\subsection*{You'll want to know: \hi{बेचैन}}
\textbf{\hi{बेचैन}} (\emph{becain}) — \textquotedblleft{}restless, uneasy\textquotedblright{}.

\textbf{\hi{बे}-} is \textquotedblleft{}without\textquotedblright{}, and \textbf{\hi{चैन}} is \textquotedblleft{}peace\textquotedblright{}: restless is \emph{without peace}.

\begin{cousinweb}
\textbf{\hi{बेचैन}} (\emph{becain}) is \textquotedblleft{}without ease\textquotedblright{}. \emph{be-} is the Persian \textquotedblleft{}without\textquotedblright{}; \emph{cain}, \textquotedblleft{}ease, rest\textquotedblright{}, is of uncertain origin, and CDIAL has no entry for it.
\end{cousinweb}

\begin{grammarlens}[title={What you've built}]
A word that shows its parts.
\end{grammarlens}

\subsection*{Your turn}
\begin{itemize}
\raggedright
  \item \emph{Say it:} \emph{becain}
  \item \emph{Say it:} \emph{becain}, once more
  \item \emph{Say it:} say \emph{maiṁ becain hūṁ}
  \item \emph{Recall:} say the Hindi for calm, then the Hindi for nervousness, then say \emph{becain} again
\end{itemize}

\begin{tcolorbox}[breakable,colback=teal!4,colframe=teal!35!black,title={Before you move on}]
What does \hi{बेचैन} mean? (\textquotedblleft{}Restless, uneasy\textquotedblright{}.) Say it once more.
\end{tcolorbox}
\section[\texorpdfstring{\hi{मन} \hi{नहीं} \hi{है} (man nahīṁ hai)}{मन नहीं है (man nahīṁ hai)}]{\hi{मन} \hi{नहीं} \hi{है} (man nahīṁ hai) — I don't feel like it}

\label{lesson:HI-C125-mannahin}



\begin{quote}
Before the new one: say the Hindi for nervousness, then the Hindi for restless.
\end{quote}

\subsection*{You'll want to know: \hi{मन} \hi{नहीं} \hi{है}}
\textbf{\hi{मन} \hi{नहीं} \hi{है}} (\emph{man nahīṁ hai}) — \textquotedblleft{}I don't feel like it\textquotedblright{}, literally \textquotedblleft{}there is no mind for it\textquotedblright{}.

\textbf{\hi{मन}} is \textquotedblleft{}mind, heart\textquotedblright{}. \textbf{\hi{आज} \hi{खाने} \hi{का} \hi{मन} \hi{नहीं} \hi{है}} — \textquotedblleft{}I don't feel like eating today\textquotedblright{}.

\begin{cousinweb}
\textbf{\hi{मन} \hi{नहीं} \hi{है}} (\emph{man nahīṁ hai}) is \textquotedblleft{}the mind is not there\textquotedblright{}. \emph{man}, \textquotedblleft{}mind, heart\textquotedblright{}, comes from Sanskrit \emph{manas}, \textquotedblleft{}mind\textquotedblright{} (CDIAL 9822).
\end{cousinweb}

\begin{grammarlens}[title={What you've built}]
A gentle refusal.
\end{grammarlens}

\subsection*{Your turn}
\begin{itemize}
\raggedright
  \item \emph{Say it:} \emph{man nahīṁ hai}
  \item \emph{Say it:} \emph{man nahīṁ hai}, once more
  \item \emph{Say it:} say \emph{āj man nahīṁ hai}
  \item \emph{Recall:} say the Hindi for nervousness, then the Hindi for restless, then say \emph{man nahīṁ hai} again
\end{itemize}

\begin{tcolorbox}[breakable,colback=teal!4,colframe=teal!35!black,title={Before you move on}]
What does \hi{मन} \hi{नहीं} \hi{है} mean? (\textquotedblleft{}I don't feel like it\textquotedblright{}.) Say it once more.
\end{tcolorbox}
\section[\texorpdfstring{\hi{मज़ा} (mazā)}{मज़ा (mazā)}]{\hi{मज़ा} (mazā) — fun, enjoyment}

\label{lesson:HI-C125-maza}



\begin{quote}
Before the new one: say the Hindi for restless, then the Hindi for I don't feel like it.
\end{quote}

\subsection*{You'll want to know: \hi{मज़ा}}
\textbf{\hi{मज़ा}} (\emph{mazā}) — \textquotedblleft{}fun, enjoyment\textquotedblright{}.

\textbf{\hi{मज़ा} \hi{आया}} — \textquotedblleft{}I enjoyed it\textquotedblright{}, literally \textquotedblleft{}fun came\textquotedblright{}. It comes to you, like laughter.

\begin{cousinweb}
\textbf{\hi{मज़ा}} (\emph{mazā}) is a loan, usually traced to Persian \emph{maza}, \textquotedblleft{}taste, relish\textquotedblright{}.
\end{cousinweb}

\begin{grammarlens}[title={What you've built}]
Calm, nervous, restless, not in the mood, and fun.
\end{grammarlens}

\subsection*{Your turn}
\begin{itemize}
\raggedright
  \item \emph{Say it:} \emph{mazā}
  \item \emph{Say it:} \emph{mazā}, once more
  \item \emph{Say it:} say \emph{mazā āyā!}
  \item \emph{Recall:} say the Hindi for restless, then the Hindi for I don't feel like it, then say \emph{mazā} again
\end{itemize}

\begin{tcolorbox}[breakable,colback=teal!4,colframe=teal!35!black,title={Before you move on}]
What does \hi{मज़ा} mean? (\textquotedblleft{}Fun, enjoyment\textquotedblright{}.) Say it once more.
\end{tcolorbox}
